\documentclass[12pt]{article}

% =========================
% Toggle: show only problems, or problems + student solution boxes
% =========================
\newif\ifsolutions
\solutionstrue     % Use this for the student template with solution boxes.
%\solutionsfalse   % POF template: questions only.

% =========================
% Packages
% =========================
\usepackage[a4paper,margin=24mm]{geometry}
\usepackage{amsmath,amssymb,mathtools}
\usepackage{braket}
\usepackage{siunitx}
\usepackage[most]{tcolorbox}
\usepackage{enumitem}
\usepackage{graphicx}
\usepackage{xcolor}
\usepackage[T1]{fontenc}
\usepackage{lmodern}
\usepackage{microtype}
\usepackage{float} % Allows figure to be placed at exact desired location
\usepackage[hidelinks]{hyperref}

\sisetup{inter-unit-product = \ensuremath{{}\cdot{}}}
\renewcommand{\d}{\mathop{}\!\mathrm{d}}
\setlength{\parindent}{0pt}
\setlength{\parskip}{0.55em}
\setlist[enumerate,1]{label=(\alph*),itemsep=0.25em,topsep=0.25em}
\setlist[enumerate,2]{label=(\roman*),itemsep=0.15em,topsep=0.15em}

% =========================
% Boxes
% =========================
\tcbset{
  enhanced,
  breakable,
  boxrule=0.8pt,
  arc=1.5mm,
  left=2mm,right=2mm,top=1.5mm,bottom=1.5mm
}

\newtcolorbox{problembox}[1]{
  colback=gray!8,
  colframe=black,
  fonttitle=\bfseries,
  title={#1}
}

\newtcolorbox{solutionbox}{
  colback=blue!3,
  colframe=blue!60!black,
  title={Solution},
  fonttitle=\bfseries
}

\newtcolorbox{answerbox}{
  colback=yellow!10,
  colframe=orange!70!black,
  title={Final answer},
  fonttitle=\bfseries
}

\newtcolorbox{insightbox}{
  colback=green!5,
  colframe=green!50!black,
  title={Physical insight},
  fonttitle=\bfseries
}

\newcommand{\problemsep}{\vspace{0.8em}}

\newcommand{\placeholderfigure}[2][50mm]{%
	\begin{center}
		\fbox{\parbox[c][#1][c]{0.82\linewidth}{\centering\small Figure placeholder}}\\[0.3em]
		{\footnotesize\emph{Caption: #2}}
\end{center}}

\graphicspath{{figures/}} % If you want to place figures in the folder 'figures' (cleaner)

\setcounter{secnumdepth}{0} % Avoids numbering of the section and subsection while allowing them to appear in the TOC



% ======================================================
% ======================================================
% ================                    ==================
% ================   BEGIN DOCUMENT   ==================
% ================                    ==================
% ======================================================
% ======================================================

\begin{document}

\begin{center}
	{\Large\bfseries QT01 -- Foundations of Quantum Mechanics\\[0.2em]
		Portfolio (POF) -- Week 8\\[0.3em]
		{\normalsize Weight: 70\% of module grade}\\[1em]
		Candidate number: 123456  % Replace candidate number with yours
	}
\end{center}


% ======================================================
% ======================================================
%                      Week 2
% ======================================================
% ======================================================
\section{Week 2}

% ======================================================
%                  Problem 1
% ======================================================
\begin{problembox}{Problem 1 \hfill [16 marks]}
At time $t=0$ a particle is represented by the wave function
\[
\psi(x,0)=
\begin{cases}
A\left(\dfrac{x}{a}\right), & 0\le x\le a,\\[0.8em]
A\dfrac{(b-x)}{(b-a)}, & a\le x\le b,\\[0.8em]
0, & \text{otherwise},
\end{cases}
\]
where $A$, $a$, and $b$ are positive real constants.

\begin{enumerate}
    \item Normalise $\psi$ (that is, find $A$, in terms of $a$ and $b$).
    \item Sketch $\psi(x,0)$, as a function of $x$.
    \item Where is the particle most likely to be found, at $t=0$?
    \item What is the probability of finding the particle to the left of $a$? Check your result in the limiting cases $b=a$ and $b=2a$.
    \item What is the expectation value of $x$?
\end{enumerate}
\end{problembox}

% ======================================================
%                  Solution 1
% ======================================================
\ifsolutions
\begin{solutionbox}
	 Replace this placeholder with your own detailed derivation, making your algebra clear and showing the steps needed to reach the result.
\end{solutionbox}

\begin{answerbox}
	State your final answer(s) clearly here.
\end{answerbox}

\begin{insightbox}
	Add a short physical interpretation here. Explain what the result means, why it matters, or what feature of the system it illustrates.
\end{insightbox}
\fi

% ======================================================
%                  Problem 2
% ======================================================
\newpage
\begin{problembox}{Problem 2 \hfill [6 marks]}
In the concrete representation of bras and kets by row and column vectors, the inner product is defined in terms of components:
\[
\braket{B|A}
=
\begin{pmatrix}
\beta_1^* & \beta_2^* & \beta_3^*
\end{pmatrix}
\begin{pmatrix}
\alpha_1\\
\alpha_2\\
\alpha_3
\end{pmatrix}
=\beta_1^*\alpha_1+\beta_2^*\alpha_2+\beta_3^*\alpha_3.
\]

The rule for inner products is essentially the same as for dot products: add the products of corresponding components of the vectors whose inner product is being calculated.

Show that the inner product $\braket{B|A}$, defined as above, satisfies all the axioms of inner products (conjugate symmetry, linearity in the second argument, and positive-definiteness).
\end{problembox}

% ======================================================
%                  Solution 2
% ======================================================
\ifsolutions
\begin{solutionbox}
	Replace this placeholder with your own detailed derivation, making your algebra clear and showing the steps needed to reach the result.
\end{solutionbox}

\begin{answerbox}
	State your final answer(s) clearly here.
\end{answerbox}

\begin{insightbox}
	Add a short physical interpretation here. Explain what the result means, why it matters, or what feature of the system it illustrates.
\end{insightbox}
\fi


\newpage
% ======================================================
% ======================================================
%                      Week 3
% ======================================================
% ======================================================
\section{Week 3}

% ======================================================
%                  Problem 1
% ======================================================
\begin{problembox}{Problem 1 \hfill [12 marks]}
Show that the position operator 
\[
\hat{x}=x,
\]
and momentum operator
\[
\hat{p}=\frac{\hbar}{i}\frac{\partial}{\partial x},
\]
are Hermitian, in other words, that
\[
\braket{\phi|\hat{x}\psi}=\braket{\hat{x}\phi|\psi}
\qquad\text{and}\qquad
\braket{\phi|\hat{p}\psi}=\braket{\hat{p}\phi|\psi}.
\]
\end{problembox}

% ======================================================
%                  Solution 1
% ======================================================
\ifsolutions
\begin{solutionbox}
	Replace this placeholder with your own detailed derivation, making your algebra clear and showing the steps needed to reach the result.
\end{solutionbox}

\begin{answerbox}
	State your final answer(s) clearly here.
\end{answerbox}

\begin{insightbox}
	Add a short physical interpretation here. Explain what the result means, why it matters, or what feature of the system it illustrates.
\end{insightbox}
\fi

% ======================================================
%                  Problem 2
% ======================================================
\newpage
\begin{problembox}{Problem 2 \hfill [8 marks]}
The wave function of an electron in the lowest energy state of a hydrogen atom is proportional to $e^{-r/a_0}$, with $a_0$ a constant and $r$ the distance from the nucleus. Calculate the relative probabilities of finding the electron inside a region of volume $1.0\,\mathrm{pm}^3$, located at
\begin{itemize}
    \item the nucleus,
    \item a distance $a_0$ from the nucleus.
\end{itemize}

As the region of interest is so small, even on the scale of the atom, we can ignore the variation of $\psi$ within it.
\end{problembox}

% ======================================================
%                  Solution 2
% ======================================================
\ifsolutions
\begin{solutionbox}
	Replace this placeholder with your own detailed derivation, making your algebra clear and showing the steps needed to reach the result.
\end{solutionbox}

\begin{answerbox}
	State your final answer(s) clearly here.
\end{answerbox}

\begin{insightbox}
	Add a short physical interpretation here. Explain what the result means, why it matters, or what feature of the system it illustrates.
\end{insightbox}
\fi


\newpage
% ======================================================
% ======================================================
%                      Week 4
% ======================================================
% ======================================================
\section{Week 4}

% ======================================================
%                  Problem 1
% ======================================================
\begin{problembox}{Problem 1 \hfill [5 marks]}
Consider a system whose state is given in terms of an orthonormal set of three vectors: $\ket{\phi_1}$, $\ket{\phi_2}$, $\ket{\phi_3}$ as
\[
\ket{\psi}=\frac{\sqrt{3}}{3}\ket{\phi_1}+\frac{2}{3}\ket{\phi_2}+\frac{\sqrt{2}}{3}\ket{\phi_3}.
\]

\begin{enumerate}
    \item Verify that $\ket{\psi}$ is normalised. Then, calculate the probability of finding the system in any one of the states $\ket{\phi_1}$, $\ket{\phi_2}$ and $\ket{\phi_3}$. Verify that the total probability is equal to one.
    \item Consider now an ensemble of 810 identical systems, each one of them in the state $\psi$. If measurements are done on all of them, how many systems will be found in each of the states $\ket{\phi_1}$, $\ket{\phi_2}$ and $\ket{\phi_3}$?
\end{enumerate}
\end{problembox}

% ======================================================
%                  Solution 1
% ======================================================
\ifsolutions
\begin{solutionbox}
	Replace this placeholder with your own detailed derivation, making your algebra clear and showing the steps needed to reach the result.
\end{solutionbox}

\begin{answerbox}
	State your final answer(s) clearly here.
\end{answerbox}

\begin{insightbox}
	Add a short physical interpretation here. Explain what the result means, why it matters, or what feature of the system it illustrates.
\end{insightbox}
\fi

% ======================================================
%                  Problem 2
% ======================================================
\newpage
\begin{problembox}{Problem 2 \hfill [17 marks]}
\begin{enumerate}
    \item Show that the sum of two Hermitian operators is Hermitian.
    \item Suppose $\hat{Q}$ is Hermitian, and $\alpha$ is a complex number. Under what condition (on $\alpha$) is $\alpha\hat{Q}$ Hermitian?
    \item When is the product of two Hermitian operators Hermitian?
    \item Show that the Hamiltonian operator
    \[
        \hat{H}=-\frac{\hbar^2}{2m}\frac{\d^2}{\d x^2}+V(x)
    \]
    is Hermitian.
\end{enumerate}
\end{problembox}

% ======================================================
%                  Solution 2
% ======================================================
\ifsolutions
\begin{solutionbox}
	Replace this placeholder with your own detailed derivation, making your algebra clear and showing the steps needed to reach the result.
\end{solutionbox}

\begin{answerbox}
	State your final answer(s) clearly here.
\end{answerbox}

\begin{insightbox}
	Add a short physical interpretation here. Explain what the result means, why it matters, or what feature of the system it illustrates.
\end{insightbox}
\fi


\newpage
% ======================================================
% ======================================================
%                      Week 5
% ======================================================
% ======================================================
\section{Week 5}

% ======================================================
%                  Problem 1
% ======================================================
\begin{problembox}{Problem 1 \hfill [20 marks]}
Consider a system whose wave function at time $t=0$ is given by
\[
\ket{\psi(x,0)}=
\frac{5}{\sqrt{50}}\ket{\phi_0(x)}
+\frac{4}{\sqrt{50}}\ket{\phi_1(x)}
+\frac{3}{\sqrt{50}}\ket{\phi_2(x)},
\]
where $\ket{\phi_n(x)}$ is the wave function of the $n^{\mathrm{th}}$ excited state for a harmonic oscillator of energy
\[
E_n=\hbar\omega\left(n+\frac{1}{2}\right).
\]

\begin{enumerate}
    \item Find the average energy of this system.
    \item Find the state $\ket{\psi(x,t)}$ at a later time and the average value of the energy; compare the result with the value obtained in (a).

    \textit{Hint: the eigenvalue equation for the time-evolution operator is}
    \[
    \hat{U}\ket{\phi_n(x)}\equiv e^{-i\hat{H}t/\hbar}\ket{\phi_n(x)}
    =e^{-iE_nt/\hbar}\ket{\phi_n(x)}.
    \]

    \item Find the expectation value of the operator $\hat{x}$ with respect to the state $\ket{\psi(x,t)}$ (i.e., find $\braket{\psi(x,t)|\hat{x}|\psi(x,t)}$).

    \textit{Hint: Use}
    \[
    \hat{x}=\sqrt{\frac{\hbar}{2m\omega}}\left(\hat{a}^{\dagger}+\hat{a}\right).
    \]
\end{enumerate}
\end{problembox}

% ======================================================
%                  Solution 1
% ======================================================
\ifsolutions
\begin{solutionbox}
	Replace this placeholder with your own detailed derivation, making your algebra clear and showing the steps needed to reach the result.
\end{solutionbox}

\begin{answerbox}
	State your final answer(s) clearly here.
\end{answerbox}

\begin{insightbox}
	Add a short physical interpretation here. Explain what the result means, why it matters, or what feature of the system it illustrates.
\end{insightbox}
\fi

% ======================================================
%                  Problem 2
% ======================================================
\newpage
\begin{problembox}{Problem 2 \hfill [6 marks]}
Calculate the expectation value of the momentum for a harmonic oscillator prepared in a state which is a superposition of two number states,
\[
\ket{\psi}=\frac{1}{\sqrt{2}}\left(\ket{\psi_1}-i\ket{\psi_2}\right).
\]
\end{problembox}

% ======================================================
%                  Solution 2
% ======================================================
\ifsolutions
\begin{solutionbox}
	Replace this placeholder with your own detailed derivation, making your algebra clear and showing the steps needed to reach the result.
\end{solutionbox}

\begin{answerbox}
	State your final answer(s) clearly here.
\end{answerbox}

\begin{insightbox}
	Add a short physical interpretation here. Explain what the result means, why it matters, or what feature of the system it illustrates.
\end{insightbox}
\fi


\newpage
% ======================================================
% ======================================================
%                      Week 7
% ======================================================
% ======================================================
\section{Week 7}

% ======================================================
%                  Problem 1
% ======================================================
\begin{problembox}{Problem 1 \hfill [10 marks]}
Consider a two-spin composite state $\ket{\psi}$ prepared such that we can write
\[
\ket{\psi}=\frac{1}{\sqrt{2}}
\left(\ket{\uparrow_A\downarrow_B}-\ket{\downarrow_A\uparrow_B}\right),
\]
where, for clarity, the subscripts $A$ and $B$ denote the subspace on which Alice and Bob will measure their spin components. First, Bob measures $y$-component of the spin with the operator $\tau_y$, then Alice measures $x$-component with the operator $\sigma_x$.

\begin{enumerate}
    \item What is the expectation value of
    \[
    \braket{\sigma_x\tau_y}=\braket{\psi|\sigma_x\tau_y|\psi}?
    \]
    \item Calculate the expectation values of
    \[
    \braket{\tau_y}=\braket{\psi|\tau_y|\psi}
    \]
    and
    \[
    \braket{\sigma_x}=\braket{\psi|\sigma_x|\psi}.
    \]
    \item What is the correlation between the two measurements
    \[
    \braket{\sigma_x\tau_y}-\braket{\sigma_x}\braket{\tau_y}?
    \]
\end{enumerate}
\end{problembox}

% ======================================================
%                  Solution 1
% ======================================================
\ifsolutions
\begin{solutionbox}
	Replace this placeholder with your own detailed derivation, making your algebra clear and showing the steps needed to reach the result.
\end{solutionbox}

\begin{answerbox}
	State your final answer(s) clearly here.
\end{answerbox}

\begin{insightbox}
	Add a short physical interpretation here. Explain what the result means, why it matters, or what feature of the system it illustrates.
\end{insightbox}
\fi

\end{document}
